% FILE: content/03_model.tex

\section{Ontological Structure of the Model}
\label{sec:model}

This section presents the formal core of the Ontology of Continua (OC).
It gives the reader the unified notation, axioms, definitions, and structural
constraints needed for the later chapters on closure, synthesis, and practical
use.

OC describes continua across multiple scientific domains through one structural
language. That language is built from axes, potentials, flows, thresholds,
boundaries, cycles, and a measure of continuumness. Each continuum is also
embedded in a surrounding meta-space. All definitions in this section are
shared at the level of declared model vocabulary. Later review chapters
will sort claims by proof status and empirical support; this section supplies
the common model language those checks depend on.

\subsection{Axiomatic foundation: Level \texorpdfstring{$K_0$}{K\_0}}

Level \(K_0\) is a purely structural substrate.
It is not a physical space; it carries no time, energy, geometry, or dynamics.
Its role is to provide the minimal conditions under which any higher-level continuum is logically possible.

\paragraph{Specification of \texorpdfstring{\(K_0\)}{K\_0}}
\(K_0\) is specified by a triple
\[
    K_0 = (S,\Delta,\mathcal{C}),
\]
where:
\begin{itemize}
    \item \(S\) is a set of states;
    \item \(\Delta : S \times S \to \mathbb{R}_{\ge 0}\) is a structural difference function;
    \item \(\mathcal{C}\) is a structural relation (or family of relations) preserving distinguishability.
\end{itemize}
There is no time parameter and no evolution operator at this level.

\paragraph{Axiom 0.1 (Difference and distinguishability).}
For all \(s_1,s_2 \in S\),
\[
    \Delta(s_1,s_2) = 0 \;\Rightarrow\; s_1 = s_2.
\]
Nonzero structural difference is the minimal condition for distinguishability.
A continuum cannot exist without at least two distinguishable states.

\paragraph{Axiom 0.2 (Nontrivial threshold \texorpdfstring{$\Theta_0$}{Theta\_0}).}
There exists \(\varepsilon > 0\) such that
\[
    \forall s_1,s_2\in S,\qquad
    s_1 \neq s_2 \Rightarrow \Delta(s_1,s_2) \ge \varepsilon.
\]
The value \(\Theta_0 = \varepsilon\) acts as a minimal structural threshold: differences smaller than \(\varepsilon\) are not resolved at the level of \(K_0\).

\paragraph{Axiom 0.3 (Logical substrate).}
\(K_0\) carries no time parameter and no dynamical operator.
It does not evolve and does not generate higher levels by itself.
It specifies only logical conditions on distinguishability and the existence of nontrivial differences.
All dynamics and all notions of energy, geometry, and causality belong to levels \(K_1\) and above.

In particular, no operator defined on \(K_0\) can increase dimension:
differences at level \(K_0\) are always expressed along the existing structural axis of distinguishability.
The dimensionality associated with \(K_0\) is fixed by the embedding meta-space \(M_0\).
Here \(M_0\) denotes the minimal meta-space in which the purely structural
substrate can be evaluated for later \(K_0\to K_1\) emergence; it is not a
physical container or a dynamical environment.

\subsection{Construction of Level \texorpdfstring{$K_1$}{K\_1}}

Level \(K_1\) is the simplest genuine continuum: it introduces time, a one-dimensional axis, and basic geometric structure.

\paragraph{Specification of \texorpdfstring{\(K_1\)}{K\_1}}
The continuum \(K_1\) is defined by:
\[
    K_1 = \big(X_1,\tau_1,\Omega_1,A_1,P_1(t),J_1(t),\Theta_1,\partial\Omega_1,C_1,k_1(t)\big),
\]
where:
\begin{itemize}
    \item \(A_1\) is a single axis (a one-dimensional coordinate);
    \item \((X_1,\tau_1)\) is a topological space obtained from the structural data of \(K_0\), typically an interval \(X_1 = (a,b)\);
    \item \(\Omega_1\) is a space of admissible configurations on \(X_1\) with appropriate regularity conditions.
          In this time-dependent analytic example,
          \[
              \Omega_1 = C^0(I,H^1(X_1,V_1)) \cap C^1(I,L^2(X_1,V_1)),
          \]
          where \(I\) is the time interval and \(V_1\) is the Hilbert space
          used by the example; a finite-dimensional vector space is the
          special case used when the local model has finitely many degrees of
          freedom.
    \item \(P_1(t)\) is an energy-like potential functional on \(\Omega_1\);
    \item \(J_1(t)\) are flows derived from the variation of \(P_1(t)\);
    \item \(\Theta_1\) encodes minimal conditions for classical stability;
    \item \(\partial\Omega_1\) is the boundary where stability thresholds saturate;
    \item \(C_1\) is the set of classical cycles (for example, periodic orbits);
    \item \(k_1(t)\) measures one-dimensional continuumness according to the general definition below.
\end{itemize}

\paragraph{Transition \texorpdfstring{$\Psi_{0\to 1}$}{Psi\_0\_to\_1}}
The passage \(K_0 \to K_1\) is generated by an operator
\[
    \Psi_{0\to 1} : (S,\Delta,\mathcal{C}) \mapsto (X_1,\tau_1,A_1,\Omega_1,\Theta_1),
\]
which constructs the skeletal \(K_1\) scaffold; the energy functional,
flow family, boundary, cycles, and continuumness scalar are specified by the
subsequent \(K_1\)-specific clauses rather than by this first transition map
alone.
The scaffold construction:
\begin{itemize}
    \item identifies a family of structurally compatible states in \(S\);
    \item equips them with an order and topology \((X_1,\tau_1)\);
    \item defines the first axis \(A_1\);
    \item constructs the admissible configuration space \(\Omega_1\);
    \item induces classical thresholds \(\Theta_1\).
\end{itemize}
This is the first instance of dimensional emergence: a continuous axis appears that cannot be represented within the purely structural substrate of \(K_0\).

\subsection{General definition of a continuum}

For any level \(K\) in the hierarchy, a continuum is defined as the tuple
\[
    K = \big(\Omega(K), A(K), P(t), J(t), \Theta(K), \partial\Omega(K), C(K), k(K,t)\big),
\]
where:
\begin{itemize}
    \item \(\Omega(K)\) is a nonempty set of admissible states;
    \item \(A(K) = \{A_1,\dots,A_n\}\) is a finite set of axes of incompatible differences;
    \item \(P(t)\) is the vector of potentials on \(\Omega(K)\);
    \item \(J(t)\) is the aggregate potential-flow vector; state updates are
          handled by the evolution operator \(E\);
    \item \(\Theta(K) = \{\Theta_i\}\) is the set of threshold functions;
    \item \(\partial\Omega(K)\) is the boundary of admissible states;
    \item \(C(K)\) is the family of structurally stable cycles;
    \item \(k(K,t)\) is the measure of continuumness.
\end{itemize}

\subsection{component-witness independence of the OC verdict interface}
\label{sec:oc13-global-verdict-minimality}

The tuple above is nonredundant only in the declared component-wise verdict-witness sense used by OC Core 1.3.3. An
equivalent axiomatisation may rename a component or derive it from other
primitive data, but it must preserve the information carried by the
corresponding OC distinction if it is to preserve the same live, dead, residue,
rebirth, domain, support, and falsifier verdicts.

\paragraph{Claim-boundary result (component-wise witness independence).}
Within the declared finite semantic verdict suite, each primitive distinction represented by the current OC release tuple has a keep/drop witness. The result does not claim unrestricted component-witness independence among all possible theories or representations. If a representation removes the distinction and does
not re-encode a functionally equivalent invariant, then there exists a pair of
admissible histories that it cannot distinguish while OC assigns different
required verdicts.

\paragraph{Proof.}
Let \(\pi_p\) be one primitive distinction: identity carrier, admissible
realisation, collapse threshold, residue relation, rebirth predicate,
\(K\)-level boundary, support class, or falsifier. For each \(\pi_p\) the
release records a separating witness pair \((H_a,H_b)\) such that \(H_a\) and
\(H_b\) agree on all other retained distinctions, but the OC verdict map
requires \(V(H_a)\neq V(H_b)\). Suppose a representation \(R\) removes
\(\pi_p\) and carries no functionally equivalent invariant. Then
\(R(H_a)=R(H_b)\). Any downstream verdict decoder \(D\) must therefore satisfy
\[
    D(R(H_a))=D(R(H_b)),
\]
contradicting the verdict-preservation requirement. Hence a representation that claims to preserve this declared verdict suite must preserve the tested distinction \(\pi_p\), or else supply an explicitly equivalent invariant. Applying the argument to the independent witness pairs for all
eight distinctions proves the component-wise lower-bound route under declared witnesses for the OC verdict class.

Consequently \(\partial\Omega\), \(J\), and \(k\) may be derived in local
models only when the derivation still preserves the boundary, flow, and
continuumness verdict information. Derivation is allowed; information loss is
not.

Every continuum is assumed to be embedded in a surrounding meta-space \(M\) such that
\[
    \Omega(K)\subseteq\Omega(M), \qquad A(K)\subseteq A(M).
\]
The meta-space provides additional admissible states and axes that can host future dimensional extensions of \(K\).

\paragraph{State space and boundary}
Thresholds are represented as functions \(f_i : \overline{\Omega(K)} \to \mathbb{R}\) with
\[
    f_i(s) \le 0 \quad\text{for all } s \in \Omega(K),
\]
and the boundary is defined as
\[
    \partial\Omega(K) = \{ s \in \overline{\Omega(K)} \mid \exists\, i: f_i(s) = 0 \}.
\]
This definition requires no metric and applies equally to geometric, topological, energetic, informational, or logical continua.

The boundary is dynamic. The notation below is an update map, not a derivative
of a set-valued object:
\[
    \partial\Omega(K,t+dt) =
    R\big(\partial\Omega(K,t),P(t),J(t),\Theta(K)\big),
\]
which can contract, expand, or bifurcate \(\Omega(K)\) during birth, life, and death events.

\subsection{Taxonomy of thresholds}

Each continuum has a structured set of thresholds \(\Theta(K)\), organised into the following types:

\begin{itemize}
    \item \textbf{Existence thresholds} \(\Theta_{\mathrm{exist}}\): conditions under which \(\Omega(K)\neq\emptyset\).
    \item \textbf{Stability thresholds} \(\Theta_{\mathrm{stab}}\): conditions for bounded dynamics and non-divergent flows.
    \item \textbf{Critical thresholds} \(\Theta_{\mathrm{crit}}\): hypersurfaces where qualitative changes (phase transitions, bifurcations) occur.
    \item \textbf{Dimensional thresholds} \(\Theta_{\mathrm{dim}}\): conditions under which new axes and higher-dimensional continua can emerge.
    \item \textbf{Death thresholds} \(\Theta_{\mathrm{death}}\): structural limits where no admissible states remain and \(\Omega(K)\) collapses.
\end{itemize}

The full threshold landscape of a continuum is thus a collection of inequalities:
\[
    f_{i}^{(\mathrm{exist})}(s) \le 0,\quad
    f_{i}^{(\mathrm{stab})}(s) \le 0,\quad
    f_{i}^{(\mathrm{crit})}(s) \le 0,\quad
    f_{i}^{(\mathrm{dim})}(s) \le 0,\quad
    f_{i}^{(\mathrm{death})}(s) \le 0,
\]
with corresponding boundary components given by the equalities.

\subsection{Potentials, flows, and structural tension}

Potentials \(P(t)\) encode the internal configuration of constraints and driving forces within a continuum.
They may correspond to energy landscapes, chemical concentrations, membrane gradients, representational or informational structures, or institutional and normative pressures.

The aggregate flow vector \(J(t)\) describes the rate of change of potentials:
\[
    \frac{dP}{dt} = J(t),
\]
with \(J(t)\) decomposed into structurally distinct components:
\begin{itemize}
    \item \textbf{supporting flows} \(J_{\mathrm{support}}\), which maintain cycles and stabilise \(k(K,t)\);
    \item \textbf{critical flows} \(J_{\mathrm{critical}}\), which move the system towards or across critical thresholds \(\Theta_{\mathrm{crit}}\) and \(\Theta_{\mathrm{dim}}\);
    \item \textbf{destructive flows} \(J_{\mathrm{kill}}\), which push the system across \(\Theta_{\mathrm{death}}\) and reduce \(k(K,t)\).
\end{itemize}

Structural tension \(T(K,t)\) is a functional of potentials, axes, and gradients (schematically \(T=T(P,A,\nabla P)\)).
It measures how strongly the current configuration stresses the threshold landscape.
Dimensional transitions occur when \(T(K,t)\) exceeds \(\Theta_{\mathrm{dim}}(K)\); collapse occurs when destructive flows combined with tension drive the system across \(\Theta_{\mathrm{death}}(K)\).

\subsection{Cycles and continuumness}

Cycles \(C(K)\) are closed trajectories in \(\Omega(K)\) that remain at a finite distance from the boundary and preserve continuumness.
At a structural level, cycles satisfy
\[
    L(C) = \int_C d\ell_K < \infty,\qquad
    S(C) = \min_{s \in C} \delta_{\partial\Omega}(s) > 0,
\]
where \(d\ell_K\) is the chosen structural length element for the realised
continuum and \(\delta_{\partial\Omega}(s)\) is the corresponding boundary
clearance.
The boundary definition above is metric-free; these cycle estimates add a
structural gauge only when quantitative clearance or length is being measured.
In metric realisations this gauge may be an induced distance. In non-metric
realisations it is any level-admissible separation functional that vanishes
exactly on the boundary and is positive on structurally interior states.

When the continuum evolves, the admissible region is read as time-dependent:
\(\Omega(K,t)\) denotes the state space available to \(K\) at time \(t\), while
\(\Omega(K)\) denotes the underlying level-relative admissible space when no
time slice is being specified.
When older formulas write \(\Omega(K(t))\), read that notation as the same
time-slice object \(\Omega(K,t)\).
The model adopts the following unified definition of continuumness:
\[
    k(K,t)
    = \chi_{\Omega}(s_K(t);K,t)\;
      S_{\mathrm{axes}}(K,t)\;
      S_{\mathrm{cycles}}(K,t)\;
      S_{\mathrm{flows}}(K,t)\;
      S_{\mathrm{coh}}(K,t),
\]
where:
\begin{itemize}
    \item \(s_K(t)\) is the current state of \(K\) at time \(t\), and
          \(\chi_{\Omega}(s_K(t);K,t)\) is the pointwise admissibility
          indicator:
          it is \(1\) if \(s_K(t)\in\Omega(K,t)\), and \(0\) otherwise;
    \item \(S_{\mathrm{axes}}(K,t)\) quantifies effective axis saturation, e.g.\ the ratio of the effective rank of working axes to the maximal possible rank for that level;
    \item \(S_{\mathrm{cycles}}(K,t)\) measures the strength and efficiency of stable cycles, such as a weighted sum of cycle efficiencies normalised by their maximal values;
    \item \(S_{\mathrm{flows}}(K,t)\) captures flow stability. With aggregated
          supporting and destructive magnitudes
          \(\Phi_{\mathrm{support}}\) and \(\Phi_{\mathrm{kill}}\), one local
          rule is shown below. The zero-flow branch applies only when the
          tested state has no active flow channel:
          \[
              S_{\mathrm{flows}}(K,t)=
              \begin{cases}
              \dfrac{\Phi_{\mathrm{support}}(t)}
                    {\Phi_{\mathrm{support}}(t)+\Phi_{\mathrm{kill}}(t)},
                    & \Phi_{\mathrm{support}}(t)+\Phi_{\mathrm{kill}}(t)>0,\\[0.7em]
              1, & \Phi_{\mathrm{support}}(t)=\Phi_{\mathrm{kill}}(t)=0.
              \end{cases}
          \]
          Thus if support vanishes while destructive flow is positive,
          \(S_{\mathrm{flows}}=0\);
    \item \(S_{\mathrm{coh}}(K,t)\in[0,1]\) encodes global structural
          coherence, which may include connectedness of the relevant graphs,
          compatibility of thresholds and flows, and absence of contradictory
          constraints.
\end{itemize}
All multiplicative factors in the displayed definition are normalised to
\([0,1]\) in the local model before the product is evaluated.
The continuumness-update operator \(U\), developed later in
the operator semantics chapter, advances \(k(K,t)\) to the next lawful time
slice using these factors.

By construction \(0 \le k(K,t) \le 1\).
Positive \(k(K,t)\) is the live-state indicator. It records four
simultaneous facts: the current state is admissible, the active axes remain
supported, the cycle factor and flow-stability factor both remain positive,
and the coherence factor has not collapsed. The full live-state status
also requires the admissible-state, threshold, and identity clauses stated in
the lifecycle and rebirth chapter.
Death corresponds to \(k(K,t)\to 0\) in combination with the collapse of \(\Omega(K)\).

\subsection{Evolution operator}

The evolution of a continuum is described at the structural level by an operator
\[
    E: K(t) \mapsto K(t+dt),
\]
or, in expanded form,
\[
    E:\ (\Omega,A,P,J,\Theta,\partial\Omega,C,k)
    \longrightarrow (\Omega',A',P',J',\Theta',\partial\Omega',C',k').
\]
Here \(A'\) denotes the post-evolution axis set of the same continuum in the
expanded tuple.
This is a unary evolution operator on a time slice of \(K\). The active
meta-space belongs to the surrounding environment. When it must be named
locally, it is shown as a subscripted parameter rather than as a second
argument of \(E\), so the operator arity remains fixed.

The model does not merge flow updates, threshold tests, cycle structure,
boundary motion, axis extension, and continuumness updates into one
differential system. It treats these components as coupled parts of a
constrained structural evolution framework with the following requirements:

\begin{itemize}
    \item \textbf{Consistency:} if \(s(t) \in \Omega(K(t))\), then either \(s(t+dt) \in \Omega(K(t+dt))\) or the transition is associated with a threshold crossing.
    \item \textbf{Threshold-respecting:} flows obey the inequality structure encoded in \(\Theta(K)\), except when explicitly crossing critical or death thresholds.
    \item \textbf{Monotonicity of dimension:} if \(A'\) contains a new axis not representable as a combination of previous axes, then \(\dim(K(t+dt)) > \dim(K(t))\); dimension cannot decrease.
\end{itemize}

The dynamics continue as long as \(\Omega(K(t)) \neq \emptyset\).
Once \(\Omega(K(t^\ast)) = \emptyset\), the continuum is dead and \(E\) can no longer act meaningfully on it.

\subsection{Embedding into meta-spaces}

Each level \(K_x\) is embedded in a meta-space \(M_x\) that provides additional admissible states and axes.
At the structural level this is captured by the conditions
\[
    \Omega(K_x(t)) \subseteq \Omega(M_x(t)), \qquad
    A(K_x) \subseteq A(M_x).
\]
The general evolution axiom can then be stated schematically as
\[
    K_x(t+dt) = E_{M_x(t)}\big(K_x(t)\big),
\]
with the additional requirement that if the time-slice embedding condition
fails, \(\Omega(K_x(t))\not\subseteq\Omega(M_x(t))\), the continuum \(K_x\)
ceases to exist. When the time slice is clear, the shorter notations \(K_x\)
and \(M_x\) suppress the argument \(t\).
Dimensional growth of \(K_x\) is only possible if the meta-space contains an
axis \(A_{\mathrm{new}}\in A(M_x)\setminus A(K_x)\) and structural tension
exceeds the corresponding dimensional threshold. This is the axis-origin
constraint used below: new axes cannot be generated from \(K_x\) alone.

\subsection{Birth of continua}

The emergence of a new continuum \(K_{x+1}\) from \(K_x\) is a threshold-induced phase transition.
Structurally, birth occurs when the following conditions are met:

\begin{enumerate}
    \item A new class of differences appears that cannot be represented using the existing axes \(A(K_x)\).
    \item Structural tension associated with these differences satisfies
          \[
              T(K_x,t) > \Theta_{\mathrm{dim}}(K_x).
          \]
    \item The embedding space \(M_x\) contains at least one axis that can host the new differences (i.e.\ \(A_{\mathrm{new}} \in A(M_x) \setminus A(K_x)\)).
    \item There exists a nonempty set of admissible states \(\Omega(K_{x+1})\) compatible with the new axis and thresholds.
\end{enumerate}

The operator of dimensional birth
\[
    \Psi_{x\to x+1}: K_x \to K_{x+1}
\]
is minimal and irreversible: any nonzero emergence of the new axis constitutes the new continuum \(K_{x+1}\), and the dimension of \(K_{x+1}\) cannot revert to that of \(K_x\) without destroying \(\Omega(K_{x+1})\).
The local model therefore instantiates dimensional monotonicity. The
corresponding no-spontaneous-creation result is recorded separately in
the spontaneous-dimension theorem discussion.

\subsection{Life of continua}

A continuum \(K\) is alive on an interval of time if
\[
    \Omega(K(t)) \neq \emptyset,\qquad
    k(K,t) > 0,\qquad
    C(K(t)) \neq \emptyset,
\]
and if supporting flows \(J_{\mathrm{support}}\) dominate destructive flows \(J_{\mathrm{kill}}\) when integrated over relevant cycles.

Life thus corresponds to the persistent existence of:
\begin{itemize}
    \item stable cycles separated from the boundary by a finite structural distance;
    \item sufficient supporting flows to maintain these cycles;
    \item coherence between potentials, axes, and thresholds;
    \item controlled critical behaviour that does not cross \(\Theta_{\mathrm{death}}\).
\end{itemize}

These conditions are interpreted differently at each level \(K_x\), but the structural pattern is the same from protocells to institutions.

\subsection{Death of continua}

A continuum \(K\) dies at time \(t^\ast\) when
\[
    \Omega(K(t^\ast)) = \emptyset,
\]
which implies \(\chi_{\Omega}(s_K(t^\ast);K,t^\ast)=0\) and hence \(k(K,t^\ast)=0\).
Equivalently:
\begin{itemize}
    \item all stable cycles vanish, \(C(K(t^\ast)) = \emptyset\);
    \item no state satisfies the threshold inequalities concurrently;
    \item any attempted continuation of dynamics would violate at least one existence threshold \(\Theta_{\mathrm{exist}}\).
\end{itemize}

\paragraph{Irreversibility of death}
Once \(\Omega(K(t^\ast)) = \emptyset\), there is no structural operator acting within the same level that can reconstruct a nonempty \(\Omega(K)\).
Any apparent resurrection would correspond to the birth of a new continuum \(K'\) with its own thresholds and state space, not the continuation of the original \(K\).
Thus death is structurally irreversible.

\subsection{Interaction of continua}

Two continua \(K_a\) and \(K_b\) interact via an interaction operator
\[
    E_{\mathrm{int}} : (K_a,K_b) \mapsto (K_a',K_b'),
\]
acting on their combined potentials, flows, thresholds, and axes.
Structurally, several regimes are distinguished:

\begin{itemize}
    \item \textbf{Competition:} flows of one continuum hinder the maintenance of cycles in the other; typically \(k_a(t)\) and \(k_b(t)\) cannot both increase indefinitely.
    \item \textbf{Parasitism:} one continuum harvests supporting flows from another, increasing its own \(k\) at the expense of the host.
    \item \textbf{Symbiosis:} supporting flows are coupled such that both continua increase their continuumness and extend their admissible regions.
    \item \textbf{Fusion:} axes and potentials combine to form a new continuum \(K_{\mathrm{fusion}}\) with merged axes and a new state space \(\Omega_{\mathrm{fusion}}\).
\end{itemize}

Interaction can itself induce dimensional transitions when mixed differences
create new, previously unused axes and drive tension above
\(\Theta_{\mathrm{dim}}\).

\subsection{Vertical hierarchy \texorpdfstring{$K_0$–$K_{12}$}{K0--K12}}

Within this formal scheme, the OC hierarchy \(K_0,\dots,K_{12}\) can be summarised as follows:

\begin{itemize}
    \item \(K_0\): purely structural substrate of distinguishable states and differences, no time, no dynamics.
    \item \(K_1\): one-dimensional classical continuum with energy functionals and basic stability thresholds.
    \item \(K_2\): physical continua, including fields, phase transitions, percolation, Berezinskii--Kosterlitz--Thouless-type transitions, and mass structure in field theory.
    \item \(K_3\): chemical continua, including reaction networks, reflexively autocatalytic and food-generated structures, concentrations, and environmental parameters.
    \item \(K_4\): protocellular continua with membranes, osmotic and curvature thresholds, internal/external gradients, and metabolic subspaces.
    \item \(K_5\): early neural and bioelectrical continua with ion channels, excitation thresholds, gradient-driven flows, and protospikes.
    \item \(K_6\): cognitive continua with representational axes, binding, internal models, memory thresholds, and prediction thresholds.
    \item \(K_7\): social continua with trust thresholds, institutional cycles, communication, and coordination flows.
    \item \(K_8\): civilisational continua with large-scale infrastructures, systemic thresholds, and long-range cycles.
    \item \(K_9\): theoretical continua comprising theories, paradigms, ontologies, and model families, with consistency and coherence thresholds.
    \item \(K_{10}\): meta-theoretical continua governing comparison, evaluation, and transformation across classes of theories.
    \item \(K_{11}\): recursive meta-theoretical continua governing operator families, meta-logical update, and model-space recursion.
    \item \(K_{12}\): global coherence continua that bound compatibility across families of \(K_{11}\)-structures.
\end{itemize}

Each level inherits the general continuum structure and adds new axes, potentials, thresholds, and cycles specific to its domain.
This chapter does not rederive every domain-specific representation theorem in
full detail here; it records the common structural backbone that those
instantiations must respect.

\subsection{Summary}

This section has presented the ontological backbone of OC:
the axioms of \(K_0\), the construction of \(K_1\), the general definition of a continuum,
the taxonomy of thresholds, the roles of potentials, flows, and structural tension,
and the definition of cycles and continuumness.
It has also fixed the evolution operator, the embedding into meta-spaces, the
structural conditions for birth, life, and death, the interaction operator, and
the vertical hierarchy \(K_0\)–\(K_{12}\).
All further developments in the Core and in extension papers rely on this structure as their common foundation.
